\documentclass[11pt, letterpaper]{article}
% ============================================================
% preamble.tex — Shared LaTeX preamble for Learning to Autolens
% ============================================================
% Contains: packages, physics notation macros, formatting.
% Include in each module via % ============================================================
% preamble.tex — Shared LaTeX preamble for Learning to Autolens
% ============================================================
% Contains: packages, physics notation macros, formatting.
% Include in each module via % ============================================================
% preamble.tex — Shared LaTeX preamble for Learning to Autolens
% ============================================================
% Contains: packages, physics notation macros, formatting.
% Include in each module via \input{../preamble}
% ============================================================

% --- Packages ---
\usepackage[utf8]{inputenc}
\usepackage[T1]{fontenc}
\usepackage{amsmath, amssymb, amsthm}
\usepackage{physics}       % \vb, \grad, \div, \curl, \dd, etc.
\usepackage{mathtools}     % \coloneqq, \prescript
\usepackage{bm}            % \bm for bold math
\usepackage{graphicx}
\usepackage{float}
\usepackage{booktabs}      % Professional tables
\usepackage{hyperref}
\usepackage{cleveref}      % \cref for smart cross-refs
\usepackage{xcolor}
\usepackage[margin=1in]{geometry}
\usepackage{enumitem}      % Custom list formatting
\usepackage{fancyhdr}      % Headers and footers
\usepackage{tcolorbox}     % Colored boxes for key results
\usepackage{listings}      % Code listings

% --- Color scheme ---
\definecolor{codeblue}{RGB}{31, 119, 180}
\definecolor{codegray}{RGB}{128, 128, 128}
\definecolor{keyresult}{RGB}{39, 123, 69}
\definecolor{exercise}{RGB}{178, 102, 0}

% --- tcolorbox environments ---
\newtcolorbox{keyresult}[1][]{
  colback=keyresult!5, colframe=keyresult!70,
  fonttitle=\bfseries, title={Key Result}, #1
}

\newtcolorbox{pythonbox}[1][]{
  colback=codeblue!5, colframe=codeblue!50,
  fonttitle=\bfseries, title={PyAutoLens Implementation}, #1
}

\newtcolorbox{exercisebox}[1][]{
  colback=exercise!5, colframe=exercise!50,
  fonttitle=\bfseries, title={Exercise}, #1
}

% --- Code listing style ---
\lstset{
  language=Python,
  basicstyle=\ttfamily\small,
  keywordstyle=\color{codeblue}\bfseries,
  commentstyle=\color{codegray}\itshape,
  stringstyle=\color{keyresult},
  breaklines=true,
  frame=single,
  framerule=0.5pt,
  rulecolor=\color{codegray!50},
}

% --- Physics macros ---
% Vectors (bold upright)
\newcommand{\vtheta}{\bm{\theta}}       % Image-plane position
\newcommand{\vbeta}{\bm{\beta}}         % Source-plane position
\newcommand{\valpha}{\bm{\alpha}}       % Deflection angle
\newcommand{\vx}{\bm{x}}               % Generic position

% Lensing quantities
\newcommand{\thetaE}{\theta_{\rm E}}    % Einstein radius
\newcommand{\SigmaCr}{\Sigma_{\rm cr}}  % Critical surface density
\newcommand{\Dd}{D_{\rm d}}             % Angular diameter distance to lens
\newcommand{\Ds}{D_{\rm s}}             % Angular diameter distance to source
\newcommand{\Dds}{D_{\rm ds}}           % Angular diameter distance lens-source

% Statistical quantities
\newcommand{\chisq}{\chi^2}             % Chi-squared
\newcommand{\chisqred}{\chi^2_{\rm red}} % Reduced chi-squared
\newcommand{\Like}{\mathcal{L}}         % Likelihood
\newcommand{\Evid}{\mathcal{Z}}         % Evidence
\newcommand{\Post}{\mathcal{P}}         % Posterior

% Abbreviations for references
\newcommand{\CK}{C\&K}                  % Congdon & Keeton (2018)
\newcommand{\NB}{N\&B}                  % Narayan & Bartelmann (1997)
\newcommand{\Night}{Nightingale+18}     % Nightingale, Dye & Massey (2018)

% --- Headers ---
\pagestyle{fancy}
\fancyhf{}
\fancyhead[L]{\small\textit{Learning to Autolens}}
\fancyhead[R]{\small\thepage}
\renewcommand{\headrulewidth}{0.4pt}

% --- Theorem environments ---
\theoremstyle{definition}
\newtheorem{definition}{Definition}[section]
\newtheorem{example}{Example}[section]

% ============================================================

% --- Packages ---
\usepackage[utf8]{inputenc}
\usepackage[T1]{fontenc}
\usepackage{amsmath, amssymb, amsthm}
\usepackage{physics}       % \vb, \grad, \div, \curl, \dd, etc.
\usepackage{mathtools}     % \coloneqq, \prescript
\usepackage{bm}            % \bm for bold math
\usepackage{graphicx}
\usepackage{float}
\usepackage{booktabs}      % Professional tables
\usepackage{hyperref}
\usepackage{cleveref}      % \cref for smart cross-refs
\usepackage{xcolor}
\usepackage[margin=1in]{geometry}
\usepackage{enumitem}      % Custom list formatting
\usepackage{fancyhdr}      % Headers and footers
\usepackage{tcolorbox}     % Colored boxes for key results
\usepackage{listings}      % Code listings

% --- Color scheme ---
\definecolor{codeblue}{RGB}{31, 119, 180}
\definecolor{codegray}{RGB}{128, 128, 128}
\definecolor{keyresult}{RGB}{39, 123, 69}
\definecolor{exercise}{RGB}{178, 102, 0}

% --- tcolorbox environments ---
\newtcolorbox{keyresult}[1][]{
  colback=keyresult!5, colframe=keyresult!70,
  fonttitle=\bfseries, title={Key Result}, #1
}

\newtcolorbox{pythonbox}[1][]{
  colback=codeblue!5, colframe=codeblue!50,
  fonttitle=\bfseries, title={PyAutoLens Implementation}, #1
}

\newtcolorbox{exercisebox}[1][]{
  colback=exercise!5, colframe=exercise!50,
  fonttitle=\bfseries, title={Exercise}, #1
}

% --- Code listing style ---
\lstset{
  language=Python,
  basicstyle=\ttfamily\small,
  keywordstyle=\color{codeblue}\bfseries,
  commentstyle=\color{codegray}\itshape,
  stringstyle=\color{keyresult},
  breaklines=true,
  frame=single,
  framerule=0.5pt,
  rulecolor=\color{codegray!50},
}

% --- Physics macros ---
% Vectors (bold upright)
\newcommand{\vtheta}{\bm{\theta}}       % Image-plane position
\newcommand{\vbeta}{\bm{\beta}}         % Source-plane position
\newcommand{\valpha}{\bm{\alpha}}       % Deflection angle
\newcommand{\vx}{\bm{x}}               % Generic position

% Lensing quantities
\newcommand{\thetaE}{\theta_{\rm E}}    % Einstein radius
\newcommand{\SigmaCr}{\Sigma_{\rm cr}}  % Critical surface density
\newcommand{\Dd}{D_{\rm d}}             % Angular diameter distance to lens
\newcommand{\Ds}{D_{\rm s}}             % Angular diameter distance to source
\newcommand{\Dds}{D_{\rm ds}}           % Angular diameter distance lens-source

% Statistical quantities
\newcommand{\chisq}{\chi^2}             % Chi-squared
\newcommand{\chisqred}{\chi^2_{\rm red}} % Reduced chi-squared
\newcommand{\Like}{\mathcal{L}}         % Likelihood
\newcommand{\Evid}{\mathcal{Z}}         % Evidence
\newcommand{\Post}{\mathcal{P}}         % Posterior

% Abbreviations for references
\newcommand{\CK}{C\&K}                  % Congdon & Keeton (2018)
\newcommand{\NB}{N\&B}                  % Narayan & Bartelmann (1997)
\newcommand{\Night}{Nightingale+18}     % Nightingale, Dye & Massey (2018)

% --- Headers ---
\pagestyle{fancy}
\fancyhf{}
\fancyhead[L]{\small\textit{Learning to Autolens}}
\fancyhead[R]{\small\thepage}
\renewcommand{\headrulewidth}{0.4pt}

% --- Theorem environments ---
\theoremstyle{definition}
\newtheorem{definition}{Definition}[section]
\newtheorem{example}{Example}[section]

% ============================================================

% --- Packages ---
\usepackage[utf8]{inputenc}
\usepackage[T1]{fontenc}
\usepackage{amsmath, amssymb, amsthm}
\usepackage{physics}       % \vb, \grad, \div, \curl, \dd, etc.
\usepackage{mathtools}     % \coloneqq, \prescript
\usepackage{bm}            % \bm for bold math
\usepackage{graphicx}
\usepackage{float}
\usepackage{booktabs}      % Professional tables
\usepackage{hyperref}
\usepackage{cleveref}      % \cref for smart cross-refs
\usepackage{xcolor}
\usepackage[margin=1in]{geometry}
\usepackage{enumitem}      % Custom list formatting
\usepackage{fancyhdr}      % Headers and footers
\usepackage{tcolorbox}     % Colored boxes for key results
\usepackage{listings}      % Code listings

% --- Color scheme ---
\definecolor{codeblue}{RGB}{31, 119, 180}
\definecolor{codegray}{RGB}{128, 128, 128}
\definecolor{keyresult}{RGB}{39, 123, 69}
\definecolor{exercise}{RGB}{178, 102, 0}

% --- tcolorbox environments ---
\newtcolorbox{keyresult}[1][]{
  colback=keyresult!5, colframe=keyresult!70,
  fonttitle=\bfseries, title={Key Result}, #1
}

\newtcolorbox{pythonbox}[1][]{
  colback=codeblue!5, colframe=codeblue!50,
  fonttitle=\bfseries, title={PyAutoLens Implementation}, #1
}

\newtcolorbox{exercisebox}[1][]{
  colback=exercise!5, colframe=exercise!50,
  fonttitle=\bfseries, title={Exercise}, #1
}

% --- Code listing style ---
\lstset{
  language=Python,
  basicstyle=\ttfamily\small,
  keywordstyle=\color{codeblue}\bfseries,
  commentstyle=\color{codegray}\itshape,
  stringstyle=\color{keyresult},
  breaklines=true,
  frame=single,
  framerule=0.5pt,
  rulecolor=\color{codegray!50},
}

% --- Physics macros ---
% Vectors (bold upright)
\newcommand{\vtheta}{\bm{\theta}}       % Image-plane position
\newcommand{\vbeta}{\bm{\beta}}         % Source-plane position
\newcommand{\valpha}{\bm{\alpha}}       % Deflection angle
\newcommand{\vx}{\bm{x}}               % Generic position

% Lensing quantities
\newcommand{\thetaE}{\theta_{\rm E}}    % Einstein radius
\newcommand{\SigmaCr}{\Sigma_{\rm cr}}  % Critical surface density
\newcommand{\Dd}{D_{\rm d}}             % Angular diameter distance to lens
\newcommand{\Ds}{D_{\rm s}}             % Angular diameter distance to source
\newcommand{\Dds}{D_{\rm ds}}           % Angular diameter distance lens-source

% Statistical quantities
\newcommand{\chisq}{\chi^2}             % Chi-squared
\newcommand{\chisqred}{\chi^2_{\rm red}} % Reduced chi-squared
\newcommand{\Like}{\mathcal{L}}         % Likelihood
\newcommand{\Evid}{\mathcal{Z}}         % Evidence
\newcommand{\Post}{\mathcal{P}}         % Posterior

% Abbreviations for references
\newcommand{\CK}{C\&K}                  % Congdon & Keeton (2018)
\newcommand{\NB}{N\&B}                  % Narayan & Bartelmann (1997)
\newcommand{\Night}{Nightingale+18}     % Nightingale, Dye & Massey (2018)

% --- Headers ---
\pagestyle{fancy}
\fancyhf{}
\fancyhead[L]{\small\textit{Learning to Autolens}}
\fancyhead[R]{\small\thepage}
\renewcommand{\headrulewidth}{0.4pt}

% --- Theorem environments ---
\theoremstyle{definition}
\newtheorem{definition}{Definition}[section]
\newtheorem{example}{Example}[section]


\fancyhead[C]{\small\textbf{Module 03: Your First Lens Model}}

\title{\textbf{Module 03: Your First Lens Model}\\[0.5em]
\large Theory Companion — Learning to Autolens}
\author{Rodrigo C\'ordova Rosado\\
\small Harvard-Smithsonian Center for Astrophysics\\
\small \texttt{rodrigo.cordova\_rosado@cfa.harvard.edu}}
\date{\today}

\begin{document}
\maketitle

\begin{abstract}
This document covers the Bayesian framework for lens modeling: the
Gaussian likelihood, prior specification, posterior inference via nested
sampling, and the key degeneracies that arise in strong lensing. This
is the theoretical foundation for the inverse problem solved
computationally in Module~03.
References: \CK\ Ch.~8; \Night\ Sec.~4--5; Skilling (2004).
\end{abstract}

\tableofcontents
\newpage

% ============================================================
\section{Bayesian Framework}
\label{sec:bayes}
% ============================================================

\subsection{Bayes' Theorem}

Given observed data $\mathbf{d}$ and model parameters $\bm{\eta}$,
Bayes' theorem gives the \textbf{posterior} probability:

\begin{keyresult}
\begin{equation}
P(\bm{\eta} \mid \mathbf{d}) = \frac{\Like(\mathbf{d} \mid \bm{\eta})\,\pi(\bm{\eta})}{\Evid}
\label{eq:bayes}
\end{equation}
\end{keyresult}

\noindent where:
\begin{itemize}[nosep]
  \item $\Like(\mathbf{d} \mid \bm{\eta})$ is the \textbf{likelihood}:
    how well parameters $\bm{\eta}$ reproduce the data
  \item $\pi(\bm{\eta})$ is the \textbf{prior}: what we know before
    seeing data
  \item $\Evid = \int \Like \, \pi \, d\bm{\eta}$ is the
    \textbf{evidence}: normalizes the posterior and is used for model
    comparison
\end{itemize}

\subsection{The Gaussian Likelihood}

For imaging data with per-pixel Gaussian noise (valid when photon
counts are large), the log-likelihood is (\Night\ eq.~3; \CK\ eq.~8.7):

\begin{keyresult}
\begin{equation}
\ln \Like = -\frac{1}{2} \sum_{i=1}^{N_{\rm pix}}
\left[ \frac{(d_i - m_i)^2}{\sigma_i^2} + \ln(2\pi\sigma_i^2) \right]
= -\frac{1}{2}\chisq + \text{const}
\label{eq:log-likelihood}
\end{equation}
\end{keyresult}

\noindent where $d_i$ is the observed pixel value, $m_i$ is the model
prediction (PSF-convolved tracer image), $\sigma_i$ is the noise in
pixel $i$, and:
\begin{equation}
\chisq = \sum_{i=1}^{N_{\rm pix}} \frac{(d_i - m_i)^2}{\sigma_i^2}
\label{eq:chi-squared}
\end{equation}

The \textbf{reduced chi-squared} is:
\begin{equation}
\chisqred = \frac{\chisq}{N_{\rm dof}}
= \frac{\chisq}{N_{\rm pix} - N_{\rm params}}
\label{eq:chi-squared-red}
\end{equation}
A good fit has $\chisqred \approx 1$. Significantly larger values
indicate an inadequate model; significantly smaller values suggest
overfitting or overestimated noise.

% ============================================================
\section{Priors}
\label{sec:priors}
% ============================================================

PyAutoLens supports three prior types:

\begin{table}[H]
\centering
\begin{tabular}{lll}
\toprule
Prior & Distribution & Use case \\
\midrule
Uniform & $\pi(\eta) = 1/(b-a)$ for $\eta \in [a,b]$ & Bounded range, no preference \\
Gaussian & $\pi(\eta) \propto \exp[-(\eta-\mu)^2/(2\sigma^2)]$ & Peaked at expected value \\
Log-Uniform & $\pi(\eta) \propto 1/\eta$ for $\eta \in [a,b]$ & Scale-invariant \\
\bottomrule
\end{tabular}
\end{table}

\textbf{Physical priors} for galaxy-scale lenses:
\begin{itemize}[nosep]
  \item $\thetaE \in [0.5'', 3.0'']$ (Uniform) — typical galaxy-scale range
  \item Lens centre: $\mathcal{N}(0, 0.1'')$ (Gaussian) — near cutout center
  \item Source position: $\mathcal{N}(0, 0.3'')$ — within caustic
  \item External shear: $|\gamma| < 0.2$ — environment contribution
\end{itemize}

% ============================================================
\section{Nested Sampling}
\label{sec:nested-sampling}
% ============================================================

\subsection{The Algorithm (Skilling 2004)}

Nested sampling computes the evidence integral by transforming the
$N$-dimensional integral into a 1D integral over \textbf{prior volume}
$X$:
\begin{equation}
\Evid = \int \Like(\bm{\eta})\,\pi(\bm{\eta})\,d\bm{\eta}
= \int_0^1 \Like(X)\,dX
\label{eq:evidence-integral}
\end{equation}
where $X(\lambda) = \int_{\Like(\bm{\eta}) > \lambda} \pi(\bm{\eta})\,d\bm{\eta}$
is the prior volume with likelihood above threshold $\lambda$.

The algorithm maintains $n_{\rm live}$ \textbf{live points} drawn from
the prior. At each iteration, the lowest-likelihood point is removed
and replaced with a new point drawn from the prior subject to
$\Like > \Like_{\rm min}$. The prior volume contracts as:
\begin{equation}
X_k \approx e^{-k/n_{\rm live}}
\label{eq:volume-contraction}
\end{equation}

\subsection{Why Nested Sampling for Lensing?}

\begin{enumerate}[nosep]
  \item Computes the evidence $\Evid$ directly — enables model comparison
  \item Explores the full prior volume — robust against multimodality
  \item Naturally produces posterior samples (as a byproduct)
  \item Handles the non-smooth likelihood surfaces common in lensing
\end{enumerate}

% ============================================================
\section{Key Degeneracies}
\label{sec:degeneracies}
% ============================================================

\subsection{Mass-Sheet Degeneracy (MSD)}

The \textbf{fundamental degeneracy} of strong lensing (\CK\ Sec.~9.1;
S92 Sec.~5.4). The transformation:
\begin{align}
\kappa'(\vtheta) &= \lambda\kappa(\vtheta) + (1-\lambda) \label{eq:msd-kappa} \\
\vbeta' &= \lambda\vbeta \label{eq:msd-beta}
\end{align}
leaves all image positions unchanged. This means lensing alone cannot
determine the absolute mass normalization.

\textbf{Breaking the MSD} requires external information:
\begin{itemize}[nosep]
  \item Time delays: $\Delta t' = \lambda\,\Delta t$ (\CK\ Sec.~9.2)
  \item Stellar kinematics: independent mass measurement
  \item Extended source morphology: constrains the source-plane magnification
\end{itemize}

\subsection{Source Size--Magnification Degeneracy}

A larger, fainter source with higher magnification can mimic a smaller,
brighter source with lower magnification. This is related to the MSD
and is partially broken by resolved source structure.

\subsection{Ellipticity--Shear Degeneracy}

Lens ellipticity and external shear produce similar quadrupole
distortions and are partially degenerate. This manifests as a banana-shaped
posterior in the $(\epsilon, \gamma)$ plane.

% ============================================================
\section{Fisher Information and Parameter Uncertainties}
\label{sec:fisher}
% ============================================================

Near the maximum likelihood, the posterior is approximately Gaussian with
covariance given by the inverse \textbf{Fisher information matrix}:
\begin{equation}
F_{ij} = -\left\langle \frac{\partial^2 \ln\Like}{\partial\eta_i \partial\eta_j} \right\rangle
\label{eq:fisher}
\end{equation}

The marginal uncertainty on parameter $\eta_i$ is:
\begin{equation}
\sigma_{\eta_i} \geq (F^{-1})_{ii}^{1/2}
\label{eq:cramer-rao}
\end{equation}
(Cram\'er-Rao bound). This gives a lower limit on achievable precision
for a given dataset and model.

% ============================================================
\section*{References}
% ============================================================

\begin{itemize}[nosep, leftmargin=*]
  \item Congdon, A.B.\ \& Keeton, C.R.\ (2018), \emph{Principles of Gravitational Lensing}, Springer
  \item Narayan, R.\ \& Bartelmann, M.\ (1997), arXiv:astro-ph/9606001
  \item Nightingale, J.W., Dye, S.\ \& Massey, R.J.\ (2018), MNRAS, 478, 4738
  \item Schneider, P., Ehlers, J.\ \& Falco, E.E.\ (1992), \emph{Gravitational Lenses}, Springer
  \item Skilling, J.\ (2004), AIP Conf.\ Proc.\ 735, 395
  \item Speagle, J.S.\ (2020), MNRAS, 493, 3132
\end{itemize}

\end{document}
